\section{Discussion and Limitations}

\paragraph{What the result establishes.}
Two streams can have the same marginal covariance spectrum and the same
expected number of raw observations, yet different learning exponents. The
relevant object is the innovation spectrum $\lambda_j/\tau_j$. A scalar
effective sample size remains adequate under uniform persistence, but
mode-dependent persistence creates a hierarchy of effective sample sizes and
moves the model--data frontier.

\paragraph{What remains open.}
First, the fixed-raw-horizon theorem is noiseless and assumes finite-variance
dwell times. Noisy stopping rules and the heavy-dwell regime $r\ge a-1$ can
exhibit additional renewal fluctuations and remain open. Second, the
coordinate model isolates coverage and is intentionally sparse. Dense Gaussian
or random feature designs may combine this mechanism with sample-covariance
concentration; proving universality beyond spectral sampling is an important
next step. Third, the clean noisy identity uses deterministic replicated
blocks or block averages. Geometric dwell lengths can add finite-block
corrections to inverse run-length moments. Finally, current experiments are
synthetic. A full empirical study should estimate mode-wise persistence on
real sequential regression datasets and test whether the innovation spectrum
predicts out-of-range learning curves better than a global correlation time.

\paragraph{Practical implication.}
The result suggests measuring not only how much variance lies in a feature
direction, but how often that direction is refreshed. When weak directions
arrive in long bursts, collecting a longer contiguous trajectory can be less
valuable than increasing temporal or environmental diversity, even when the
number of logged observations is unchanged.

